% MSc dissertation — Actor-Aware and Robust Measurement of Bitcoin Mining Decentralisation
% Built on the School of Informatics infthesis skeleton (structure preserved).
%
% Leave one of the documentclass lines uncommented to match your degree.
% You may remove the logo option if it causes problems.
% Do not change any other options.
% \documentclass[logo,msc,adi]{infthesis}     % Adv Design Inf
% \documentclass[logo,msc,ai]{infthesis}      % AI
% \documentclass[logo,msc,cogsci]{infthesis}  % Cognitive Sci
% \documentclass[logo,msc,cs]{infthesis}      % Computer Sci
% \documentclass[logo,msc,cyber]{infthesis}   % Cyber Sec
% \documentclass[logo,msc,datasci]{infthesis} % Data Sci
% \documentclass[logo,msc,di]{infthesis}      % Design Inf
% \documentclass[logo,msc,dsti]{infthesis}    % Data Sci TI
% \documentclass[logo,msc,inf]{infthesis}     % Informatics
\documentclass[logo,msc]{infthesis}           % degree unspecified, do not change except to add your degree
%%%%%%%%%%%%%%%%%%%%%%%%
% Understand any problems and seek approval before assuming it's ok to remove ugcheck.
\usepackage{msccheck}

% Include any packages you need below, but don't include any that change the page
% layout or style of the dissertation.
\usepackage{microtype} % recommended, but you can remove if it causes problems
% --- Added (layout-safe; verify the document still compiles under msccheck) ---
\usepackage{graphicx}  % for \includegraphics (your result figures)
\usepackage{booktabs}  % for clean tables (\toprule/\midrule/\bottomrule)
\usepackage{amsmath}   % for the metric definitions / equations
% Put your PNGs in a subfolder ./figures/ ; uncomment to set the path:
 \graphicspath{{figures/}}

\begin{document}
\begin{preliminary}

\title{Actor-Aware and Robust Measurement of Bitcoin Mining Decentralisation}

\author{Evgenios Pilia Kyriakou}

\date{\today}

\abstract{
Claims about the (de)centralisation of Bitcoin mining rest on attributing each block
to a mining pool, an uncertain step most studies treat as settled. This dissertation
treats block attribution as the object of study and measures how sensitive
decentralisation conclusions are to the choices it involves. Over the full Bitcoin
block history, we compute six decentralisation metrics under a systematic grid of
attribution choices, actor definitions, tag/address heuristics, and unknown-block
treatments, reporting uncertainty as sensitivity bands rather than point estimates.
Bitcoin mining follows a U-shaped trajectory, decentralising to its most decentralised
state around 2017 and re-concentrating over the years since, and these turning points
are robust where coverage supports measurement. We establish an ordering of uncertainty: how blocks are attributed matters
far more than sampling, and within attribution the tag-versus-address heuristic
dominates. We further show that the two major public reference lists---which share community
lineage---now agree to near-identical attribution, relocating attribution uncertainty
away from list choice and onto the tag-versus-address split, coverage, and unknown handling;
that operator clustering via shared payout addresses is confounded; and recommend that
decentralisation studies report sensitivity envelopes and confidence intervals with
explicit disclosure of attribution method and coverage.
}

\maketitle

\newenvironment{ethics}
   {\begin{frontenv}{Research Ethics Approval}{\LARGE}}
   {\end{frontenv}\newpage}

\begin{ethics}
This project was planned in accordance with the Informatics Research
Ethics policy. It did not involve any aspects that required approval
from the Informatics Research Ethics committee.

\standarddeclaration
\end{ethics}


\begin{acknowledgements}
% TODO(you): not yet written. Thank your supervisor (Dr Daniel W. Woods) and
% anyone else here, or delete the environment entirely before submission.
\end{acknowledgements}


\tableofcontents
\end{preliminary}


%======================================================================
\chapter{Introduction}
\label{ch:intro}
% Drop-in replacement for the Introduction chapter in skeleton.tex.
% BUDGET ~4 pp. Faithful LaTeX rendering of the supplied draft, with
% citations wired to mybibfile.bib and cross-references to later chapters.
% NOTE: \S\ref{sec:intro-structure} kept consistent with the actual
% chapter layout of this thesis (Intro / Background / Methodology /
% Results / Discussion / Conclusions) -- see flag in chat.
 
\section{Background and motivation}
\label{sec:intro-motivation}
 
The degree to which Bitcoin is decentralised at its consensus layer is among
the most consequential properties one can ask about the system. It underpins
security assessments --- a small coalition of block producers can, in
principle, censor transactions or attempt to rewrite recent history --- and it
increasingly carries legal and financial weight, informing the ``sufficiently
decentralised'' tests that several securities regulators invoke when
classifying a digital asset, as well as institutional investment and risk
decisions. In practice this abstract property is operationalised through a
handful of concentration statistics computed over mining entities: the
Nakamoto coefficient, the Herfindahl--Hirschman Index (HHI), the Gini
coefficient, Shannon entropy, and the concentration ratios $\mathrm{CR}_3$ and
$\mathrm{CR}_5$. Each summarises, in a different way, how block production is
distributed across the entities that produce blocks.
 
Computing any of these statistics, however, presupposes an answer to a prior
and surprisingly difficult question: which real-world entity produced each
block? Because the protocol does not record the identity of a block's
producer, this must be inferred, almost always by matching a
signature-like string in the block's coinbase transaction (a ``tag''), or the
coinbase payout address, against curated lists of known mining pools. This
coinbase-based methodology is standard and pervasive, yet the sensitivity of
the resulting decentralisation estimates to the many choices it involves has
never been characterised systematically. The metric sits at the end of a
measurement pipeline; the attribution step sits at its very beginning; and the
literature has scrutinised the former far more closely than the latter.
 
\section{Problem statement}
\label{sec:intro-problem}
 
Two practical difficulties follow from treating attribution as a solved
preliminary. The first is \emph{comparability}. Published decentralisation
estimates may differ substantially depending on which heuristic is used, how
conflicting signals from the tag and the address are resolved, and how blocks
that match neither are handled, and these choices are rarely disclosed. Two
studies reporting different Nakamoto coefficients for the same period may
simply have made different, undocumented attribution decisions. The second is
\emph{reliability}. Even an internally consistent methodology may yield
measurements that are unstable at the daily or weekly granularities at which
decentralisation is often reported, so that apparent movements reflect
estimation noise rather than genuine structural change.
 
A further complication, rarely acknowledged, is that the entity signing a
block's coinbase transaction is not necessarily the entity that selected the
block's contents. The growth of external block-template providers,
payout-routing services, and Stratum~v2 job negotiation means that ``the pool
that mined a block'' and ``the entity that chose what went into it'' are
increasingly distinct. On-chain data alone cannot resolve this separation;
this dissertation therefore fixes its attribution target explicitly as the
\emph{signing pool}, and treats the template-provider layer as a known
limitation of any coinbase-based approach rather than a problem it claims to
solve.
 
Across the existing literature, attribution uncertainty is either ignored,
documented without being propagated, or silently absorbed into a single
best-guess series. Romiti et al.~\cite{romiti2019deep}, combining several
attribution sources over roughly 556{,}400 blocks, found 684 blocks that
different sources assigned to different entities, and even blockchain.com
labels its own hashrate charts as rough estimates, direct evidence that
attribution is an active, unresolved problem. Yet no prior study, to my
knowledge, has (i)~formally propagated attribution uncertainty through to the
final decentralisation metrics, (ii)~decomposed how much of the downstream
variation each specific upstream choice is responsible for, or (iii)~tested
whether the identification of structural change in decentralisation over time
is itself an artefact of methodology. These three gaps define the space this
dissertation occupies.
 
\section{Aims and objectives}
\label{sec:intro-aims}
 
This dissertation reframes the measurement task. Rather than producing yet
another estimate of pool shares, it asks whether consensus-layer
decentralisation estimates are stable, well-specified, and robust to the
attribution choices that precede them. It builds directly on the measurement
framework of Ovezik et al.~\cite{ovezik2025sok}, which systematised the
decentralisation-measurement workflow and demonstrated that downstream
pre-processing choices materially affect the result, but stopped short of the
upstream attribution stage; and on the earlier stratified framework of
Karakostas et al.~\cite{ovezik2022sok}, which identified actor identification
as a central unresolved challenge. The two bodies of work are complementary,
and together they motivate carrying sensitivity analysis all the way back to
the step at which blocks are first mapped to entities. The concrete objectives
are:
 
\begin{description}
  \item[O1.] build a reproducible, end-to-end attribution pipeline over public
    Bitcoin data;
  \item[O2.] quantify the sensitivity of the standard decentralisation metrics
    to three families of attribution choice, actor definition, attribution
    heuristic, and the treatment of unattributed blocks;
  \item[O3.] characterise the temporal granularity at which the metrics are
    statistically stable; and
  \item[O4.] derive principled reporting recommendations for future
    decentralisation studies and for the Edinburgh Decentralisation Index.
\end{description}
 
\section{Research questions}
\label{sec:intro-rqs}
 
Four research questions operationalise these objectives:
 
\begin{description}
  \item[RQ1 - Actor-definition sensitivity.] To what extent do
    decentralisation metrics vary under alternative definitions of the relevant
    actors?
  \item[RQ2 - Workflow sensitivity.] How much do the metrics change under
    different attribution heuristics and different treatments of unattributed
    blocks?
  \item[RQ3 - Measurement stability.] At what temporal granularity (daily,
    weekly, monthly, quarterly) do the metrics achieve sufficient stability for
    confident interpretation?
  \item[RQ4 - Change-point robustness.] Is the identification of structural
    change in decentralisation over time robust across attribution scenarios,
    or is it itself methodology-dependent?
\end{description}
 
\section{Contributions and principal findings}
\label{sec:intro-contributions}
 
The substantive picture that emerges is a U-shaped decentralisation
trajectory: Bitcoin mining reached its most decentralised state around 2017
and has been re-concentrating since, with a detected structural break in late
2022 and the present day sitting near a Nakamoto coefficient of two. Against that backdrop, the methodological
contributions of this dissertation are:
 
\begin{description}
  \item[A reproducible pipeline and attributed dataset.] An open, end-to-end
    pipeline over public Bitcoin data, and a block-level dataset carrying
    method and confidence flags, suitable for integration into the Edinburgh
    Decentralisation Index (O1).
  \item[Sensitivity envelopes and a clean ordering of uncertainty.] Replacing
    point estimates with sensitivity bands reveals the central empirical
    result: the sources of uncertainty fall into a strict order. For HHI in the
    modern era, the band induced by the choice of attribution heuristic
    ($\approx 0.27$) exceeds that from unknown-block handling ($\approx
    0.037$), which exceeds that from actor definition ($\approx 0.015$), which
    in turn exceeds the bootstrap sampling half-width ($\approx 0.004$).
    Attribution uncertainty therefore dominates sampling uncertainty by one to
    two orders of magnitude: how blocks are attributed matters far more than
    how many are sampled (O2). That the actor definition ranks third of the
    four sources is itself a result rather than a null finding --- an
    actor-aware treatment was needed to establish, rather than assume, that the
    choice is second-order for security-relevant metrics, and to isolate the
    tail-sensitive Gini as the exception.
  \item[A characterisation of temporal stability.] Under pre-registered
    criteria (a coefficient-of-variation threshold and an adjacent-window
    confidence-interval-overlap test), monthly aggregation is stable for every
    metric except the Gini coefficient, which is fragile throughout (O3, RQ3).
  \item[A test of structural-change robustness.] The recent
    re-concentration is a genuine structural shift that survives changes of
    attribution scenario, with residual instability localised precisely to the
    address-based signal; the exact dating of gradual transitions, by contrast,
    remains attribution-sensitive, a finding in its own right (RQ4).
  \item[Convergence of public reference lists.] The two major public
    pool-reference lists now disagree on only two blocks in all of Bitcoin's
    history, far below the hundreds of conflicts reported by older
    multi-source studies, so, contrary to what that earlier work might
    suggest, list choice is not where attribution uncertainty resides today; it
    lives in the tag-versus-address split, coverage, and unknown handling.
  \item[Reporting guidance.] A short set of recommendations following from the
    above: decentralisation studies should report sensitivity envelopes
    alongside confidence intervals, and disclose the attribution heuristic and
    coverage, rather than a bare point estimate (O4).
\end{description}
 
\section{Structure of the dissertation}
\label{sec:intro-structure}
 
The remainder of this dissertation is organised as follows.
Chapter~\ref{ch:background} reviews the relevant literature, the
decentralisation-measurement frameworks of Ovezik et al.\ and Karakostas et
al., the empirical attribution work of Romiti et al., and earlier
network-layer studies, and locates the three gaps this work addresses.
Chapter~\ref{ch:method} describes the methodology: the three-stage pipeline,
the decentralisation metrics, and the experimental design that varies
attribution along three orthogonal dimensions, together with the change-point
and bootstrap procedures and the pre-registered evaluation criteria.
Chapter~\ref{ch:results} reports the empirical result, attribution
coverage, the decentralisation trajectory, the sensitivity analyses answering
RQ1 and RQ2, the temporal-stability analysis (RQ3), the change-point
robustness analysis (RQ4), and the inter-list coverage findings, and
assembles them into the ordering of uncertainty sources.
Chapter~\ref{ch:discussion} distils the reporting guidance and reflects on the
limitations of the work, chiefly the signing-pool framing and the confounded
operator-clustering signal. Chapter~\ref{ch:conclusions} concludes and outlines
the probabilistic-attribution and temporal-graph extensions left as future
work.


%======================================================================
\chapter{Background and Related Work}
\label{ch:background}
% Drop-in replacement for the Background chapter in skeleton.tex.
% BUDGET ~7 pp. Faithful LaTeX rendering of the supplied draft:
% mining/pools -> attribution -> metrics -> related work + three-gap handoff.
% Metric formulas set as proper maths; citations wired to mybibfile.bib.
% NOTE: cross-references in 2.4 point to the consolidated chapter layout of
% this thesis (see flag in chat re: single Results chapter).
 
This chapter establishes the technical and scholarly context for the
dissertation. It proceeds from the mechanics of Bitcoin mining and the pools
that dominate block production (\S\ref{sec:bg-mining}), through the problem of
attributing individual blocks to those pools (\S\ref{sec:bg-attribution}), to
the metrics by which the resulting concentration is quantified
(\S\ref{sec:bg-metrics}). Section~\ref{sec:bg-related} then situates the work
against the existing literature and identifies the specific gaps it addresses.
 
\section{Bitcoin mining and mining pools}
\label{sec:bg-mining}
 
Bitcoin extends its ledger through proof of work \cite{nakamoto2008bitcoin}.
Miners assemble candidate blocks and repeatedly hash the block header, varying
a nonce, until the resulting digest falls below a network-wide difficulty
target; the first miner to succeed broadcasts the block, and the rest of the
network extends it. The target is adjusted every 2{,}016 blocks so that blocks
arrive on average roughly every ten minutes regardless of how much hashing
power is online. Because a valid digest can only be found by brute force, the
probability that any given miner produces the next block is, over the long run,
proportional to its share of the total hash rate.
 
Each block's first transaction is the \emph{coinbase transaction}, which mints
new coins and pays the block reward --- the fixed subsidy (currently 3.125~BTC,
following the 2024 halving, and itself halving every 210{,}000 blocks) plus the
transaction fees --- to one or more output addresses of the miner's choosing.
The coinbase transaction also contains an input field, the coinbase script,
that references no prior output and is therefore almost unconstrained: miners
may write arbitrary bytes into it. These two degrees of freedom, the payout
address and the free-form coinbase script, are, as
\S\ref{sec:bg-attribution} explains, exactly what makes block attribution
possible.
 
Solo mining is a high-variance activity. A miner controlling a small fraction
of the network hash rate may, on expectation, wait years between blocks, even
though its long-run share of rewards is fixed. \emph{Mining pools} exist to
remove that variance. A pool aggregates the hash rate of many participants,
directs them to work on a common block template, and distributes the reward
among them in proportion to the work each contributes, measured through
low-difficulty ``shares'', under payout schemes such as pay-per-share (PPS)
or pay-per-last-$N$-shares (PPLNS). This turns an individually unpredictable
lottery into a steady income stream.
 
The consequence is structurally important for this dissertation: even when the
underlying mining hardware is owned by a very large and geographically
dispersed population, the \emph{power to build and propose blocks} concentrates
into a small number of pool operators. Consensus-layer decentralisation is
therefore measured over pools rather than over individual machines: the pool
is the relevant actor because it, not the individual participant, decides what
a block contains and where its reward is paid.
 
A caveat qualifies even this framing. Historically the pool operator controlled
the block template, so the entity signing the coinbase and the entity choosing
the block's contents were the same. More recently, external block-template
providers, payout-routing services, and the Stratum~v2 protocol (which permits
participants to negotiate or construct their own transaction sets) have begun
to separate these roles, so that ``the pool that signed a block'' and ``the
entity that chose what went into it'' are no longer necessarily identical.
On-chain data alone cannot disentangle the two. This dissertation therefore
fixes its attribution target explicitly as the \emph{signing pool}, and treats
the template-provider layer as an acknowledged limitation of any coinbase-based
methodology rather than a problem it claims to resolve.
 
Finally, two conventions make pools identifiable at all. Because the coinbase
script is free-form, pool operators customarily embed a short human-readable
marker in it, for example \texttt{/Foundry USA/} or \texttt{/F2Pool/} , 
and they tend to pay rewards to a stable set of reused addresses. Neither
convention is enforced by the protocol; both are voluntary, inconsistent across
pools, and change over time. They are nonetheless the entire basis on which
attribution rests.
 
\section{Attributing blocks to pools}
\label{sec:bg-attribution}
 
A block does not declare who produced it. The protocol records the coinbase
script and the payout addresses, but not the identity of the miner.
\emph{Attribution} is the inference from these on-chain artefacts to a
real-world entity, performed by matching them against community-maintained
reference lists of known pools. Two independent signals are available.
 
\subsection{Two attribution signals}
 
\textbf{Coinbase-tag matching} decodes the raw bytes of the coinbase script to
text and searches for known pool marker strings. For instance, a script
containing the bytes that decode to \texttt{\dots/Foundry USA/\dots} is
attributed to Foundry USA. \textbf{Payout-address matching} instead compares
the coinbase output address against lists of addresses known to belong to
particular pools. The two signals catch overlapping but not identical sets of
blocks: a block may carry a recognisable tag but pay to an unlisted address, or
vice versa.
 
\subsection{Reference lists}
 
Both signals depend on curated reference lists that map marker strings and
addresses to canonical pool names. The most widely used are the
blockchain-etl / btc.com known-pools list and the mempool.space pool list.
Address-based identification builds on a long line of transaction-graph
clustering work, beginning with the shared-authority and change-address
heuristics of Meiklejohn et al.~\cite{meiklejohn2013fistful}, which group
addresses under common control; the reference lists used here can be seen as a
curated, pool-specific instance of the same idea.
These lists are living artefacts: pools rebrand, new pools appear, and payout
addresses rotate, so a list is only ever a snapshot of a moving target. Because
different lists are compiled independently, they can also disagree with one
another, a source of uncertainty that this dissertation makes an explicit
object of study rather than a hidden assumption.
 
\subsection{Failure modes}
 
Three things can go wrong, and each corresponds to a discretionary choice that
published estimates rarely disclose:
 
\begin{itemize}
  \item \textbf{Missing signal.} Neither the tag nor the address matches any
    list entry, leaving the block \emph{unattributed} (``unknown''). This is
    especially common in Bitcoin's early years, before tagging conventions were
    widely adopted; across the full history in this project's data, roughly
    30\% of blocks are unknown. How this pile is treated, counted as a single
    anonymous entity, excluded, or assigned to the largest known pool,
    measurably changes the result.
  \item \textbf{Conflict.} The tag names one pool while the address names
    another, and a resolution rule must decide which signal wins. Romiti et
    al.~\cite{romiti2019deep} documented 684 such direct conflicts across
    roughly 556{,}400 blocks when combining multiple attribution sources,
    concrete evidence that the choice is not merely hypothetical.
  \item \textbf{Inter-list disagreement.} Two reference lists map the same
    signal to different names, or one list recognises a pool the other omits.
    Whether this materially affects modern estimates is itself an empirical
    question that this dissertation answers.
\end{itemize}
 
\subsection{The confidence scheme}
 
Rather than force a single verdict, the two signals are combined into a
confidence scheme that records, for each block, whether tag and address agree
(\emph{high} confidence), conflict (\emph{ambiguous}), or fail to match at all
(\emph{unknown}). This exposes attribution quality as a first-class output
instead of concealing it inside a point estimate. The choices left open by this
scheme, which signal to trust, how to resolve conflicts, and what to do with
unknowns, are precisely the dimensions along which this dissertation varies
the pipeline, defined formally in Chapter~\ref{ch:method}.
 
\section{Measuring decentralisation}
\label{sec:bg-metrics}
 
Once every block in a time window has been attributed, the window yields a
distribution of block counts across entities, which normalises to a vector of
shares $s_1, s_2, \dots$ summing to one. A decentralisation metric compresses
that distribution into a single number capturing how concentrated block
production is. This dissertation computes the metrics most common in the
literature; they are summarised below. A useful property is that the
concentration-oriented and decentralisation-oriented measures move in opposite
directions, so agreement between the two families on the direction of a change
is a built-in consistency check.
 
\paragraph{Nakamoto coefficient.}
Introduced by Srinivasan and Lee~\cite{srinivasan2017quantifying}, the Nakamoto
coefficient is the minimum number of entities whose combined share exceeds
50\% --- the smallest coalition that could, in principle, seize majority
control of block production:
\[
  \mathrm{NC} = \min\Big\{k : \textstyle\sum_{i=1}^{k} s_i > \tfrac{1}{2}\Big\}.
\]
A higher value means greater decentralisation, and the metric is directly tied
to the 51\% security threshold. Being integer-valued, it is coarse and can be
brittle near boundaries, a fragility this dissertation observes empirically.
 
\paragraph{Herfindahl--Hirschman Index (HHI).}
The HHI is the sum of squared shares,
$\mathrm{HHI} = \sum_{i} s_i^2$,
ranging from $1/n$ (perfectly even) to $1$ (a monopoly), and is conventionally
scaled by $10{,}000$ to match its use in competition economics. Because shares
are squared, it is dominated by the largest entities and responds sensitively
to changes across the whole distribution. Higher values indicate greater
concentration.
 
\paragraph{Gini coefficient.}
Borrowed from the study of inequality \cite{gini1912variabilita}, the Gini
coefficient measures how unequal the share distribution is, ranging from $0$
(perfect equality) toward $1$. It captures the overall shape of the
distribution, but, as this dissertation finds, it is the most fragile of
the metrics, being especially sensitive to how the long tail of small pools and
the pile of unknown blocks are handled.
 
\paragraph{Shannon entropy.}
The Shannon entropy of the share distribution,
$\mathrm{H} = -\sum_{i} s_i \log_2 s_i$,
measured in bits and ranging from $0$ to $\log_2 n$, quantifies the
unpredictability of which entity mined a given block. Higher entropy means a
more even distribution and hence greater decentralisation.
 
\paragraph{Concentration ratios $\mathrm{CR}_3$ and $\mathrm{CR}_5$.}
The concentration ratios are simply the combined share of the top three
($\mathrm{CR}_3 = \sum_{i=1}^{3} s_i$) and top five
($\mathrm{CR}_5 = \sum_{i=1}^{5} s_i$) entities. They are intuitive and
frequently invoked in regulatory settings, with higher values indicating more
concentration. Their main limitation is saturation: once a handful of entities
dominate, they become insensitive to further shifts among the leaders.
 
Crucially, these metrics are only the final step of a longer measurement
workflow. Ovezik et al.~\cite{ovezik2025sok} systematised that workflow,
data collection, clustering and attribution, population estimation, and
windowing, and showed that the downstream choices (window size, inclusion
thresholds, population estimation) can significantly move the result. This
dissertation extends the same sensitivity lens one stage further upstream, to
the attribution step that feeds every one of these metrics.
 
\section{Related work}
\label{sec:bg-related}
 
\subsection{Frameworks for measuring decentralisation}
 
Karakostas, Kiayias and Ovezik~\cite{ovezik2022sok} proposed a stratified,
multi-layer view of blockchain decentralisation, spanning hardware,
software, network, consensus, tokenomics, client API, governance, and geography,
and, for each layer, identified the relevant resources and the parties that
control them. Critically for this project, they flagged actor identification
and clustering as a central, unresolved challenge in any empirical
decentralisation study. Their successor work, Ovezik et
al.~\cite{ovezik2025sok}, operationalised the full consensus-layer measurement
workflow and demonstrated that pre-processing choices materially affect
decentralisation estimates. Its most consequential gap, from this
dissertation's standpoint, is that the systematic sensitivity analysis it
performs for downstream choices is not carried through to the upstream
attribution stage itself, the step at which blocks are first mapped to
entities. This dissertation takes that gap as its direct starting point.
 
\subsection{Empirical evidence for attribution uncertainty}
 
Romiti et al.~\cite{romiti2019deep}, in ``A Deep Dive into Bitcoin Mining
Pools'', provide the strongest empirical grounding for the problem. Combining
several attribution sources across roughly 556{,}400 blocks, they identified
684 direct conflicts in which different sources attributed the same block to
different mining entities, and their combined methodology attributed on the
order of 35{,}000 more blocks than blockchain.com alone, which itself labels
its hashrate charts as rough estimates, explicitly conceding that some blocks
cannot be reliably attributed. These findings confirm that attribution is an
active, unsolved problem. Their limitation, for present purposes, is that they
stop at documenting inter-source disagreement and do not propagate that
uncertainty into any downstream decentralisation metric.
 
\subsection{Earlier empirical decentralisation studies}
 
Gencer et al.~\cite{gencer2018decentralization} measured network-layer
decentralisation in Bitcoin and Ethereum, and like most work in the area
proceeded with point-estimate reporting and opaque attribution. Closest in
spirit to the present work, Lin et al.~\cite{lin2021measuring} computed the
Gini coefficient, Shannon entropy and Nakamoto coefficient for Bitcoin mining
across daily, weekly and monthly granularities --- anticipating the
metric-and-granularity structure of RQ3 --- but took block-to-entity
attribution as a given input and reported point estimates rather than
propagating attribution uncertainty. At the theoretical end, Kwon et
al.~\cite{kwon2019impossibility} proved that full decentralisation is
unattainable in permissionless proof-of-work systems without a Sybil cost,
formalising \emph{why} mining power concentrates into pools in the first place
and motivating the empirical question of \emph{how much}. Together these
reinforce the generality of the concern: the reliance on undisclosed
attribution choices and single best-guess numbers is not an isolated oversight
but a pervasive feature of the literature.
 
\subsection{The gap this dissertation addresses}
 
The consistent thread across this body of work is that attribution uncertainty
is either ignored, documented without being operationalised, or silently
absorbed into a single best-guess series. To the author's knowledge, no prior
study has (i)~formally propagated attribution uncertainty through to the final
decentralisation metrics, (ii)~decomposed how much of the downstream variation
is due to each specific upstream choice, or (iii)~examined whether the
identification of structural change in decentralisation over time is itself
methodology-dependent. These three gaps define the contribution space of this
dissertation, and the chapters that follow address them in turn:
Chapter~\ref{ch:method} sets out the methodology that makes the attribution
choices explicit and measurable, and Chapter~\ref{ch:results} reports the
sensitivity, stability, and change-point results that fill each gap.

%======================================================================
\chapter{Methodology}
\label{ch:method}
% Drop-in replacement for the Methodology chapter in skeleton.tex.
% BUDGET ~9 pp. The design-space section (\ref{sec:designspace}) carries the
% supervisor-flagged argument: every choice justified by (a) exhausting the logic,
% (b) spanning the space, or (c) matching named literature positions.
 
This chapter describes the measurement pipeline and defends its design. The
standard the design must meet is stated up front: a sensitivity analysis is only
as meaningful as the set of variants it sweeps. Choose variants arbitrarily and
the resulting ``sensitivity band'' is an artifact of the menu. Every choice in
this chapter is therefore justified against one of three explicit criteria,
the variants \emph{exhaust the logical possibilities} of a choice, they
\emph{span the space} of reasonable options including its deliberate extremes,
or they \emph{correspond to named positions} in the literature, and in each
case the excluded alternatives are identified and the exclusion defended.
 
\section{Pipeline overview}
\label{sec:pipeline}
 
The analysis is organised as a three-stage pipeline, built around a single
asymmetry: expensive steps run once and are cached; cheap steps are re-run
under systematic variation.
 
\textbf{Stage 1 (extraction)} pulls every Bitcoin block from the public
BigQuery dataset \texttt{crypto\_bitcoin} and caches the fields attribution
requires, height, timestamp, the coinbase parameter (the free-form field in
which pools embed identifying tags), and the coinbase output addresses (to
which the block reward is paid), into a single columnar artifact of
954{,}375 blocks. This stage runs once for the entire project, so no analysis
depends on repeated queries against a mutable external source.
 
\textbf{Stage 2 (attribution)} maps each block to a pool label using two
independent signals, coinbase-tag matching and payout-address matching,
combined under an explicit confidence scheme (\S\ref{sec:data}). Attribution is
where every methodological choice this thesis studies is made, and it is the
stage the experimental design varies: each configuration of the design space
(\S\ref{sec:designspace}) produces one cached attribution artifact.
 
\textbf{Stage 3 (metrics)} aggregates attributed blocks into time windows and
computes six decentralisation metrics (\S\ref{sec:metrics}) per window. This
stage is pure arithmetic over cached attributions and completes in minutes,
which is what makes the full experimental grid, and the bootstrap
resampling of \S\ref{sec:uncertainty}, computationally feasible.
 
The consequence of this architecture is that the experimental design is a
Cartesian product rather than a set of bespoke analyses: heuristics, actor
definitions, and unknown-block treatments are implemented as interchangeable
functions, so any combination is produced by composition. The attribution and
metric modules carry a unit-test suite, and the entire pipeline reproduces from
a single script; code, derived data, and the pre-registered analysis
thresholds (\S\ref{sec:uncertainty}) are publicly archived
(Appendix~\ref{app:repro}).% TODO(you): ensure Appendix~\ref{app:repro} carries a real, resolvable repository/DataShare link before submission.
 
\section{Data and reference lists}
\label{sec:data}
 
\textbf{Block data.} The primary source is BigQuery's public
\texttt{crypto\_bitcoin} dataset, which mirrors the blockchain and exposes
every block's coinbase transaction. Operating a full Bitcoin node was
considered and rejected: the secondary source reliably supplies exactly the
fields attribution requires, and the marginal integrity benefit does not
justify the engineering cost at this scale. Coinbase parameters are decoded as
UTF-8 throughout, several major pools embed non-ASCII tags, and
extraction queries filter on the dataset's month-partition column to bound
cost (details in Appendix~\ref{app:supp}).
 
\textbf{Attribution target.} The target is the \emph{signing pool}, the
entity that constructed and published the block, not the individual hashers
who contribute work to it. This is the standard target in the empirical
literature \cite{romiti2019deep,gencer2018decentralization} and the only one
recoverable from on-chain data. Two limitations follow and are acknowledged
now: a pool is itself an aggregate whose internal composition is invisible
on-chain, and emerging template-negotiation protocols (Stratum~v2) may
decouple block construction from pool identity in ways that will degrade
tag-based attribution in the future (revisited in
Chapter~\ref{ch:discussion}).
 
\textbf{The two signals.} Pools voluntarily embed identifying strings (e.g.\
\texttt{/F2Pool/}) in the coinbase parameter; tag matching searches this field
for known substrings. Pools also reuse payout addresses; address matching
compares coinbase outputs against known pool addresses. The signals are
complementary in kind: a tag is a \emph{declarative} assertion by the block
producer, precise when present, but absent for non-tagging pools, while
an address is \emph{behavioural} evidence that does not depend on the pool's
cooperation but is diluted by address rotation, custodial sharing, and
rebrand inheritance. That the two signals have different failure modes is
precisely what makes the rule for combining them a substantive methodological
choice rather than an implementation detail.
 
\textbf{Reference lists.} Two independently maintained public lists supply the
tag and address references: the btc.com-lineage \texttt{Blockchain-Known-Pools}
list (149 tags, 190 addresses) as primary, and mempool.space's
\texttt{pools-v2.json} (201 tags, 218 addresses) as the second source. The
relationship between the lists is treated as an \emph{output} of the study
rather than a hidden assumption: Chapter~\ref{ch:results} measures their
disagreement at reference and block level, following Romiti et al.'s finding
that inter-source attribution conflict has historically been non-trivial
\cite{romiti2019deep}.
 
\textbf{Confidence scheme.} For each block, each signal independently returns
a label or no match; the cross-product yields five recorded outcomes,
matched by both in agreement, tag only, address only, matched by both in
conflict, and unknown. Every heuristic in \S\ref{sec:designspace} is then a
deterministic \emph{reading} of this five-way column rather than a
re-computation, guaranteeing that heuristics differ only in the combination
rule, never in the underlying matching.
 
\section{The attribution design space}
\label{sec:designspace}
 
The experimental design varies the pipeline along three orthogonal dimensions.
These are not arbitrary: they correspond to the three places where mapping raw
blocks to a share distribution forces a decision --- \emph{what counts as one
actor} (D1), \emph{which signal names the producer} (D2), and \emph{what to do
with blocks no signal names} (D3). Together they parameterise the attribution
step that the SoK measurement framework of Ovezik et al.\ identifies but holds
fixed \cite{ovezik2025sok}: their sensitivity analysis systematically varies
the downstream pre-processing of a share distribution, taking
block-to-entity attribution as a given input. Making that input a varied
dimension is the gap this thesis fills. Table~\ref{tab:designspace} summarises
the space; the subsections defend each dimension in turn.
 
\begin{table}[t]
  \centering
  \caption{The three-dimensional attribution design space. The baseline
  configuration (A1/H1/U2) is the literature's implicit default.}
  \label{tab:designspace}
  \begin{tabular}{llp{8.2cm}}
    \toprule
    Dim. & Code & Variant \\
    \midrule
    \multirow{3}{*}{D1 Actor}
      & A1 & Each distinct pool label is one actor (literature standard) \\
      & A2 & Labels below a minimum all-time share collapsed into one long-tail actor (threshold swept) \\
      & A3 & Labels inferred to share an operator are merged (graph clustering) \\
    \midrule
    \multirow{5}{*}{D2 Heuristic}
      & H1 & Tag only \\
      & H2 & Address only \\
      & H3 & Union of both signals; tag preferred on conflict \\
      & H4 & Intersection --- both signals must agree \\
      & H5 & Union of both signals; address preferred on conflict \\
    \midrule
    \multirow{3}{*}{D3 Unknowns}
      & U1 & All unattributed blocks as a single anonymous entity \\
      & U2 & Exclude unattributed blocks and renormalise \\
      & U3 & Worst case: assign all unattributed blocks to the largest known entity \\
    \bottomrule
  \end{tabular}
\end{table}
 
\subsection{Actor definitions (D1): three named literature positions}
 
Metrics are computed over a share distribution across \emph{actors}, and
``actor'' is not given by the data: a pool label is not guaranteed to be one
real-world entity, and many labels are economically negligible. The three
definitions swept are not inventions of this project but articulated positions
in the measurement literature. \textbf{A1}, one label, one actor, is the
naive default: every study that reports ``pool $X$ has share $s$'' without
further processing is using it. \textbf{A2}, collapsing the sub-threshold
long tail, formalises the ubiquitous ``top-$N$ plus Other'' reporting
convention; rather than fixing one threshold, the threshold is \emph{swept}, so
that A2 is a parametric family and the sensitivity attributable to the
threshold itself is measured rather than assumed. \textbf{A3}, merging
labels that share a real-world operator, is precisely the entity-clustering
step (2b) that the SoK framework names as the central unresolved challenge of
actor identification \cite{ovezik2025sok}. A3 is operationalised by the most
natural available method: a graph over pool labels linked by shared payout
addresses, clustered by connected components (as a sanity check) and Louvain
community detection (as the main method) \cite{blondel2008louvain}.
 
Two aspects of D1's treatment require explicit defence. First, A2 and A3 are
reported \emph{separately} throughout, because they instantiate different
steps of the SoK workflow, population inclusion versus entity clustering
and pooling them into one ``actor sensitivity'' number would silently mix
two distinct methodological questions. Second, the A3 arm is reported as a
sensitivity \emph{upper bound} rather than a measurement: as
Chapters~\ref{ch:results} and \ref{ch:discussion} show, shared-address
clustering systematically over-merges independent major pools because major
pools genuinely share large numbers of payout addresses through custodial and
mining-as-a-service infrastructure. This is itself a methodological finding,
the natural method for the SoK's open operator question was tested and
shown, with evidence, to be confounded, and it is why the clean answer to
RQ1 rests on A1 and A2.
 
\emph{Excluded alternatives.} Full transaction-graph entity resolution
(clustering the entire address graph rather than pool labels) was rejected at
proposal stage because it layers a hard-to-evaluate inference on top of the
uncertainty under study; the A3 finding vindicates that caution empirically,
since the shared-infrastructure confound identified here afflicts the full
graph a fortiori. Manual curation of operator identities from off-chain
evidence is sound but unscalable and unfalsifiable at this scope; it is the
recommended input to future probabilistic attribution
(Chapter~\ref{ch:conclusions}). Interview-based validation with pool operators
was rejected on proportionality grounds: it would introduce human-subjects
ethics scope disproportionate to the methodological benefit at MSc scale.
 
\subsection{Attribution heuristics (D2): exhausting the combination space}
 
The confidence scheme of \S\ref{sec:data} reduces attribution to two binary
events per block: the tag signal matched or did not, and the address signal
matched or did not. Any deterministic heuristic over these signals must answer
three questions: which signals does it consult; does one match suffice, or are
both required; and when both match but disagree, which wins? Enumerating the
answers yields exactly the five variants of Table~\ref{tab:designspace}. H1
and H2 are the two single-signal rules. H4 is the conjunctive rule,
attribute only when both signals agree, the precision-oriented
high-confidence tier. The disjunctive rule --- attribute when either signal
matches --- is incomplete without a conflict resolution, and there are exactly
two orderings: tag-preferred (H3) and address-preferred (H5). H1--H5 are
therefore not a sample of plausible heuristics but a \emph{complete
enumeration of the deterministic combination rules over two binary signals}.
This is the strongest available answer to ``why these and not others'': within
the deterministic setting, there are no others.
 
The enumeration is not merely formal, each cell is a position actually
occupied in practice. Tag-only attribution is the de-facto standard of block
explorers and most empirical studies; address-only is what remains when pools
do not tag; and the intersection and conflict-ordering rules matter
empirically precisely because Romiti et al.\ documented that the two signal
families disagree in the wild, with their resolution materially changing the
attributed set \cite{romiti2019deep}.
 
\emph{Excluded alternatives.} Probabilistic attribution, assigning each
block a distribution over pools rather than a label, is the natural
generalisation beyond determinism. It was deliberately deferred: a
probabilistic model introduces an additional, hard-to-validate layer of
modelling uncertainty \emph{on top of} the uncertainty this thesis measures,
and would confound the question; it is specified as future work in
Chapter~\ref{ch:conclusions}. Heuristics drawing on auxiliary evidence
(transaction-graph structure, template fingerprints, network-layer data) were
excluded because they change the \emph{evidence base} rather than the
\emph{combination rule}; the graph-based member of that family was tested
separately as A3, with the instructive outcome described above.
 
\subsection{Unknown-block treatments (D3): spanning the space}
 
Every heuristic leaves a residue of blocks no signal can name, around 30\%
of all history at baseline, far more in the early years. What is done with
this mass is a population-estimation decision: the cleanest instance in this
pipeline of the SoK framework's step 2c, which Ovezik et al.\ name as
consequential but which no prior study has propagated through to the final
metrics \cite{ovezik2025sok}. Three treatments are swept, and the
justification is a spanning argument with a specific structure. \textbf{U1},
all unknowns as one anonymous entity, and \textbf{U2}, exclude and
renormalise, are the \emph{two defensible interpretive endpoints}: treat
the unknown mass as a distinct thing, or condition it away. U2 deserves
particular attention because it is the literature's silent default: every
published pie chart of ``known pools'' is U2, and it embeds the untested
assumption that unknown blocks are distributed like known ones. \textbf{U3},
assign all unknowns to the largest known entity, is not a belief about
the world but a \emph{deliberate adversarial bound}: the most concentrated
state consistent with the observed evidence, directly relevant to the security
reading of the Nakamoto coefficient.
 
\emph{Excluded alternatives.} Any scheme that distributes unknowns
proportionally across known pools is bracketed between U2 (its limit) and
U1/U3 (its extremes), and was excluded as adding parameters without extending
the envelope, the envelope, not any interior point, is the object of
interest. The best-case mirror of U3 (assign all unknowns to a new smallest
entity) coincides with U1 for every metric except the Nakamoto coefficient and
was excluded as redundant.
 
\subsection{The design as a whole}
 
The dimensions are orthogonal by construction, any $(A,H,U)$ triple is a
valid configuration. The full heuristic-by-unknown grid (fifteen scenarios) is
materialised for the workflow-sensitivity question, while single-dimension
sweeps around a principled baseline answer the per-dimension questions. The
baseline, A1/H1/U2: label actors, tag attribution, exclude-and-renormalise
is chosen not because it is believed correct but because it is what the
literature implicitly does, which is exactly the point of departure a
sensitivity study should measure from. Finally, the results validate the menu
retrospectively: Chapter~\ref{ch:results} shows each dimension matters for at
least one metric, and the \emph{ordering} of their effects, the central
result, is something only a design that varies all of them on the same data
could deliver.
 
\section{Decentralisation metrics}
\label{sec:metrics}
 
Let a time window contain $N$ attributed blocks across $n$ actors, with actor
$i$ producing $c_i$ blocks and share $p_i = c_i/N$, ordered $p_1 \ge p_2 \ge
\dots \ge p_n$. Six metrics are computed per window, chosen to cover the
distinct facets of concentration used in the literature
\cite{srinivasan2017quantifying,gencer2018decentralization}:
 
\begin{align}
  \mathrm{NC} &= \min\Big\{k : \sum\nolimits_{i=1}^{k} p_i > \tfrac{1}{2}\Big\},
  &
  \mathrm{HHI} &= \sum\nolimits_{i=1}^{n} p_i^2,
  \notag\\
  \mathrm{H} &= -\sum\nolimits_{i=1}^{n} p_i \log_2 p_i,
  &
  \mathrm{CR}_k &= \sum\nolimits_{i=1}^{k} p_i \quad (k = 3, 5),
  \label{eq:metricdefs}
\end{align}
together with the Gini coefficient
$G = \sum_{i}\sum_{j} |p_i - p_j| \,/\, (2n \sum_i p_i)$.
 
The Nakamoto coefficient \cite{srinivasan2017quantifying} is the consensus
security threshold, the number of actors needed to assemble a strict mining
majority ($>50\%$ throughout). HHI, the standard antitrust concentration
measure, and Shannon entropy are shape measures with useful effective-number
readings: $1/\mathrm{HHI}$ and $2^{\mathrm{H}}$ both estimate the ``effective
number of pools'', from opposite emphases. $\mathrm{CR}_3$ and $\mathrm{CR}_5$
are the intuitive top-$k$ shares of regulatory practice. Gini measures
whole-distribution inequality and is maximally sensitive to the long tail.
Higher HHI, Gini and $\mathrm{CR}_k$ indicate more concentration; higher NC
and entropy indicate less. The metrics deliberately differ in what they
weight: NC and $\mathrm{CR}_k$ depend only on the head of the distribution,
HHI and entropy on its overall shape, and Gini heavily on its tail
differences the results exploit, since a choice that moves only the tail
(actor definition) should, and does, register in Gini but not NC. Monthly
windows are the baseline granularity; the choice is tested rather than
assumed (RQ3, \S\ref{sec:uncertainty}).
 
\section{Change-point method}
\label{sec:changepoint}
 
Structural breaks in each metric series are detected by exact least-squares
segmentation: dynamic programming over an $L_2$ cost, the objective family
underlying PELT \cite{killick2012pelt}. In the taxonomy of Truong et
al.~\cite{truong2020selective}, this fixes the cost function (an $L_2$
piecewise-constant-mean model) and the search method (exact dynamic
programming) while constraining the number of changes, the three components
their survey identifies for any offline change-point method. The
implementation is pure NumPy 
the standard change-point libraries do not build on the Python environment
used, and was validated against exhaustive search, agreeing on 300 of 300
randomised test cases. Because a single penalty value is not comparable across
series of very different lengths, the number of detected breaks is fixed per
series rather than sharing one penalty across granularities; this
operationalisation is disclosed because it is a choice.
 
A methodological point shapes how RQ4 is evaluated. Much of the underlying
signal is smooth drift rather than sharp shift, and on a smooth signal the
detected break \emph{date} is intrinsically unstable: minute perturbations of
the series move the optimal segmentation boundary by months without changing
the qualitative story. Demanding exact date agreement across attribution
scenarios would therefore be a brittle test that confounds attribution
sensitivity with segmentation fragility. RQ4 instead uses two complementary
measures: \emph{trajectory agreement} (rank correlation of each scenario's
full metric series against the baseline), which tests whether scenarios tell
the same story without privileging any single date; and \emph{sharp-break
corroboration} (the proportion of scenarios detecting the salient break within
a fixed tolerance window, compared against the agreement rate expected by
chance). The scope of the robustness claim is delimited in advance by the
coverage analysis: break comparison is performed on the high-coverage era that
RQ2 and RQ3 independently establish as measurable, with an unrestricted
all-history run reported as the scope check that justifies the restriction.
 
\section{Uncertainty quantification}
\label{sec:uncertainty}
 
The thesis distinguishes two orthogonal sources of uncertainty and measures
both on the same data so they can be compared.
 
\textbf{Attribution uncertainty} --- not knowing who produced each block ---
is operationalised by the design-space sweep: at each window, the spread of a
metric across configurations is its \emph{sensitivity envelope}, decomposable
by dimension (the D2 band across H1--H5; the D3 band across U1--U3; the D1
band across actor definitions). Methodological choices are thereby treated as
an uncertainty source to be propagated, not a footnote to be disclosed.
 
\textbf{Sampling uncertainty} --- even under perfect attribution, a finite
window of block production is a noisy draw from the underlying hashrate
distribution --- follows the framing of Milad et al., who develop hypothesis
tests for the Nakamoto coefficient under stochastic block arrival but take
attribution as given \cite{milad2024statistical}. Here it is operationalised
non-parametrically so that it covers all six metrics uniformly: within each
window, blocks are resampled and the metric recomputed 1{,}000 times, and the
2.5th--97.5th percentiles form a 95\% confidence interval. Two resampling
models are run as a cross-check on model dependence: \emph{multinomial}
resampling, which conditions on the window's block total and is the exact
block-level bootstrap for count-based shares, and \emph{Poisson} resampling,
which draws independent per-actor counts matching the block-arrival process.
Agreement between the two establishes that the intervals reflect sampling
variability rather than a resampling-model artifact. Intervals are computed
for \emph{every} configuration of the grid, so the two uncertainty sources are
directly comparable window by window and metric by metric, the comparison
that yields the ordering result of Chapter~\ref{ch:results}.
 
\textbf{Pre-registration.} The temporal-stability question (RQ3) is evaluated
against thresholds fixed \emph{before} the stability analysis was run and
committed to the repository with a timestamp: a window-to-window
coefficient-of-variation ceiling, a rank-correlation floor of $\rho > 0.8$
against the monthly baseline, change-point agreement within $\pm 2$ months
across granularities, and the proposal's adjacent-window bootstrap-CI-overlap
criterion ($>50\%$). Pre-registration is reported as method, not decoration:
because the thresholds predate the data, passing them tests a hypothesis
rather than decorating one.

%======================================================================
\chapter{Results}
\label{ch:results}
% Drop-in replacement for the Results chapter in skeleton.tex.
% BUDGET ~11 pp. Ordered as a narrative: coverage -> trajectory -> sensitivity
% (RQ2 -> RQ1 -> LR4 -> HR2/RQ3) -> RQ4 capstone -> synthesis.
 
This chapter presents the empirical results. The order is deliberate: we first
establish what the attribution pipeline produces and where its coverage is strong
(\S\ref{sec:res-coverage}), then present the headline decentralisation trajectory
under the baseline configuration (\S\ref{sec:res-trajectory}). The remainder of the
chapter interrogates how much of that picture survives methodological variation:
the attribution heuristic (\S\ref{sec:res-rq2}), the actor definition
(\S\ref{sec:res-rq1}), the choice of reference list (\S\ref{sec:res-interlist}),
and sampling noise (\S\ref{sec:res-sampling}), before testing whether the
trajectory's structural breaks themselves are robust (\S\ref{sec:res-rq4}).
Section~\ref{sec:res-synthesis} assembles these into the central result: a
quantitative ordering of every uncertainty source in the measurement.
 
%----------------------------------------------------------------------
\section{Attribution and coverage}
\label{sec:res-coverage}
 
The full-history pipeline processed 954{,}375 blocks. Under the baseline tag-only
attribution (H1), 70.0\% of all blocks are attributed to one of 109 distinct pool
labels, leaving 30.0\% unknown across the whole of Bitcoin's history.
Table~\ref{tab:toppools} lists the largest attributed producers.
 
\begin{table}[t]
  \centering
  \caption{Largest block producers over the full history under tag-only
  attribution (H1). Blocks are counts of attributed blocks, 2009--2026.}
  \label{tab:toppools}
  \begin{tabular}{lr@{\qquad}lr}
    \toprule
    Pool & Blocks & Pool & Blocks \\
    \midrule
    AntPool     & 107{,}908 & CloverPool (BTC.com) & 37{,}796 \\
    F2Pool      &  96{,}600 & BTC Guild            & 32{,}935 \\
    Foundry USA &  71{,}063 & Poolin               & 26{,}015 \\
    ViaBTC      &  51{,}596 & Binance Pool         & 21{,}897 \\
    Braiins     &  41{,}874 & \\
    \bottomrule
  \end{tabular}
\end{table}
 
Running both signals over the same blocks decomposes the population by evidence
type: 388{,}341 blocks are matched by tag only, 278{,}922 by both signals in
agreement (the high-confidence tier), 59{,}283 by address only, and 226{,}683 by
neither. Overall, tag matching covers 70.0\% of blocks, address matching only
35.6\%, and their union 76.2\%. Two observations from this decomposition drive
much of what follows. First, the two signals are far from interchangeable: the
address signal alone recovers barely half the coverage of the tag signal, so
heuristics that lean on it are working from a much smaller evidence base. Second,
within-list tag/address disagreement is negligible (1{,}146 blocks, and
\S\ref{sec:res-interlist} shows even this is largely an artifact), so the
substantive question is not \emph{which signal wins when they conflict} but
\emph{what happens to the blocks only one signal can reach}.
 
Figure~\ref{fig:coverage} shows the unknown-block rate over time. Coverage is
poor before roughly 2014, pooled mining was nascent and tagging conventions
had not stabilised, then high and stable in the modern era. This coverage
profile delimits the rest of the chapter: it is the reason results are reported
with a high-coverage (2018 onwards) reference period, and it foreshadows the
chapter's organising finding, namely that where coverage is low, every
downstream conclusion becomes attribution-dependent.
 
\begin{figure}[t]
  \centering
  \includegraphics[width=0.8\linewidth]{fullhistory_coverage.png}
  \caption{Attribution coverage over Bitcoin's history: the unknown-block rate
  by month under the baseline attribution. Coverage is structurally poor before
  $\sim$2014 and high thereafter.}
  \label{fig:coverage}
\end{figure}
 
%----------------------------------------------------------------------
\section{The decentralisation trajectory}
\label{sec:res-trajectory}
 
Figure~\ref{fig:trajectory} shows the six metrics over monthly windows under the
baseline configuration (tag attribution, unknowns excluded and renormalised,
one label per actor). The picture is a pronounced U-shape: Bitcoin mining
\emph{decentralised} through its first eight years, reaching its most
decentralised state around 2017, the Nakamoto coefficient peaks at 4--5,
meaning four to five independent pools were needed to assemble a mining
majority, and has been \emph{re-concentrating} since, with the onset marked by
a detected structural break in late 2022 (\S\ref{sec:res-rq4}). At the end
of the measured period the baseline sits near a Nakamoto coefficient of 2, with
the top three pools producing about 70\% of blocks ($\mathrm{CR}_3 \approx
0.70$) and the top five about 90\% ($\mathrm{CR}_5 \approx 0.90$).
 
Two change-points recur throughout the analyses that follow: a break in late
2014, as the early chaotic era gives way to stable industrial pooled mining, and
a break in late 2022, marking the onset of the recent re-concentration. Both
dates are detected identically at every temporal granularity tested
(\S\ref{sec:res-sampling}). Values before $\sim$2014 are shown for completeness
but should not be over-read: this is precisely the low-coverage era flagged in
\S\ref{sec:res-coverage}, and \S\ref{sec:res-rq4} confirms that no stable break
structure exists there.
 
This U-shape is the headline empirical picture of the thesis. Everything that
follows asks a single question of it: how much of this picture is a fact about
Bitcoin, and how much is an artifact of the choices made to draw it?
 
\begin{figure}[t]
  \centering
  \includegraphics[width=0.8\linewidth]{metrics_over_time.png}
  \caption{Decentralisation metrics over monthly windows under the baseline
  attribution (H1/U2). Mining decentralises to a $\sim$2017 peak and
  re-concentrates thereafter.}
  \label{fig:trajectory}
\end{figure}
 
%----------------------------------------------------------------------
\section{Sensitivity to attribution method (RQ2)}
\label{sec:res-rq2}
 
RQ2 asks how much the measured picture moves when the attribution workflow
changes. Two nested sweeps answer it: first varying only the unknown-block
treatment (U1--U3) at the baseline heuristic, then the full grid of heuristic
$\times$ unknown treatment (H1--H5 $\times$ U1--U3, fifteen scenarios).
Sensitivity is summarised as the \emph{band width}: at each monthly window, the
spread between the maximum and minimum metric value across scenarios, averaged
over the high-coverage era (2018 onwards).
 
\textbf{Unknown-block treatment alone moves the metrics modestly.} Across
U1--U3 the mean 2018+ band widths are 0.037 for HHI, 0.040 for Gini, 0.053 for
$\mathrm{CR}_5$ and 0.068 for $\mathrm{CR}_3$, visible, but small relative to
the range each metric traverses over the period. The U-shaped trajectory and
both turning points survive under every treatment. The exception is the Nakamoto
coefficient, whose band reaches a full unit: as a threshold statistic it changes
discontinuously when the unknown mass is re-assigned around the 50\% boundary, a
brittleness that recurs throughout this chapter.
 
\textbf{The attribution heuristic dominates.} Adding the heuristic axis widens
the envelope by a factor of three to six (Table~\ref{tab:bands};
Figure~\ref{fig:sens-heur}): the 2018+ band across H1--H5 is 0.272 for HHI
against 0.037 for the unknown axis, 1.375 against 0.233 for Shannon entropy, and
0.250 against 0.068 for $\mathrm{CR}_3$. The mechanism is visible in the
per-heuristic coverage: H1, H3 and H5 attribute 70.0--76.2\% of blocks, while H2
(address-only) and H4 (intersection) cover only 35.6\% and 29.2\%. The three
high-coverage heuristics trace an essentially identical trajectory; the wide
envelope is driven almost entirely by the two low-coverage ones, which compute
shares over a small, unrepresentative attributed subset.
 
\begin{table}[t]
  \centering
  \caption{RQ2 sensitivity band widths (mean max--min across scenarios per
  monthly window, 2018 onwards), and per-heuristic attribution coverage. The
  heuristic axis dominates the unknown-treatment axis by a factor of 3--6.}
  \label{tab:bands}
  \begin{tabular}{lccc}
    \toprule
    & Unknown axis (U1--U3) & Heuristic axis (H1--H5) & Ratio \\
    \midrule
    HHI              & 0.037 & 0.272 & $7.4\times$ \\
    Shannon entropy  & 0.233 & 1.375 & $5.9\times$ \\
    $\mathrm{CR}_3$  & 0.068 & 0.250 & $3.7\times$ \\
    \midrule
    \multicolumn{4}{l}{\emph{Coverage by heuristic:} H1 70.0\% \(\cdot\) H2 35.6\% \(\cdot\) H3 76.2\% \(\cdot\) H4 29.2\% \(\cdot\) H5 76.2\%}\\
    \bottomrule
  \end{tabular}
\end{table}
 
\begin{figure}[t]
  \centering
  \includegraphics[width=\linewidth]{sensitivity_by_heuristic.png}
  \caption{Metric trajectories by attribution heuristic. The high-coverage
  heuristics (H1/H3/H5) coincide; the envelope is driven by the low-coverage
  H2 and H4.}
  \label{fig:sens-heur}
\end{figure}
 
The finding, stated plainly: \emph{which signals a study uses to attribute
blocks matters several times more than what it does with the blocks it cannot
attribute}, and the divergence is a coverage phenomenon. This has a direct
reporting implication, developed in Chapter~\ref{ch:discussion}: a
decentralisation estimate is uninterpretable without its attribution coverage,
because coverage is the single best predictor of where in the sensitivity
envelope the estimate lies.
 
%----------------------------------------------------------------------
\section{Sensitivity to actor definition (RQ1)}
\label{sec:res-rq1}
 
RQ1 asks whether it matters what counts as \emph{one actor}. The clean answer
comes from the A2 global-threshold sweep, in which labels below a minimum
all-time share are collapsed into a single long-tail actor and the threshold is
varied across the plausible range (Figure~\ref{fig:actor}). The underlying
distribution explains the result before the metrics do: only 16 of the 109
pools ever exceed 1\% of all-time block production, so the long tail is
numerous but nearly massless.
 
Consequently the head-weighted metrics barely respond: across the full
threshold sweep the Nakamoto coefficient, HHI, $\mathrm{CR}_3$ and
$\mathrm{CR}_5$ move by 4--9\% of their values, an HHI band of roughly 0.015
over the high-coverage era, the figure carried into the uncertainty ordering of
\S\ref{sec:res-synthesis}. The Gini coefficient is the
outlier, moving by 37\%: as a whole-distribution inequality measure it is
constructed from exactly the tail mass the threshold removes. The practical
reading is symmetric, actor definition is a second-order concern for
security-oriented metrics, and a first-order disclosure requirement for anyone
reporting Gini.
 
\begin{figure}[t]
  \centering
  \includegraphics[width=\linewidth]{actor_threshold_sensitivity.png}
  \caption{Actor-definition sensitivity under the A2 threshold sweep. Head-weighted
  metrics (Nakamoto, HHI, CR) move by 4--9\%; the tail-sensitive Gini moves by
  37\%.}
  \label{fig:actor}
\end{figure}
 
The third actor definition, A3 (operator grouping via shared payout
addresses), produced apparent shifts an order of magnitude larger, an HHI
displacement of 0.266 and a Nakamoto reduction of 1.7. These numbers are
reported \emph{only as an upper bound}, not as a measurement: the clustering
that produces them systematically fuses independent major pools (AntPool,
F2Pool, ViaBTC, Binance and others into a single ``operator''), because major
pools genuinely share hundreds to thousands of payout addresses through
custodial and mining-as-a-service infrastructure. The confound is in the signal,
not the tuning, no edge-weight threshold breaks the spurious component while
preserving genuine merges, and Chapter~\ref{ch:discussion} presents the
evidence in full. RQ1's answer therefore rests on A1 and A2: actor definition
is a real but localised sensitivity, concentrated in tail-sensitive metrics.
 
%----------------------------------------------------------------------
\section{Inter-list agreement (LR4)}
\label{sec:res-interlist}
 
A natural suspicion at this point is that the results depend on \emph{whose}
reference list the pipeline uses. They do not, and establishing this
relocates the uncertainty precisely.
 
At the reference level, the two major public lists (btc.com lineage and
mempool.space) share 142 tags and 183 addresses, and after resolving a single
rebrand alias --- btc.com's ``CloverPool'' is mempool.space's ``BTC.com'' ---
there are \emph{zero} genuine conflicts between them. At the block level the
agreement is equally stark: of the 668{,}319 blocks matched by both lists,
99.9997\% receive the same attribution, two conflicting blocks in the whole
of Bitcoin's history with near-identical coverage (70.0\% vs.\ 71.4\%;
Figure~\ref{fig:interlist}). The 1{,}146-block ``conflict'' count of
\S\ref{sec:res-coverage} is thus confirmed as a within-list tag/address
artifact, not a disagreement between sources.
 
\begin{figure}[t]
  \centering
  \includegraphics[width=0.85\linewidth]{interlist_conflict.png}
  \caption{Per-list coverage and genuine inter-list conflict over time for the
  btc.com-lineage and mempool.space reference lists. The modern lists have
  converged to near-perfect agreement.}
  \label{fig:interlist}
\end{figure}
 
This is an apparent reversal of the historical picture: Romiti et
al.~\cite{romiti2019deep} found 684 attribution conflicts across the more
numerous, older and heterogeneous sources of 2019. The comparison is not
strictly like-for-like, however, and one caveat sets the correct reading:
the two lists share community lineage --- mempool.space's list descends from
the same open-source effort --- so their near-perfect agreement measures the
\emph{consolidation of the public attribution ecosystem} into a single de-facto
standard, not the independent concordance of two separately-constructed
instruments. Read that way it is still consequential for practice:
attribution uncertainty is real, but \emph{it does not live in list choice}. It
lives in the tag-versus-address split, in coverage, and in unknown handling,
exactly the dimensions of \S\ref{sec:res-rq2}.
 
%----------------------------------------------------------------------
\section{Sampling uncertainty and temporal stability (HR2, RQ3)}
\label{sec:res-sampling}
 
The uncertainty sources so far are all attributional. The remaining candidate
is statistical: even under perfect attribution, a finite window of blocks is a
noisy draw from the underlying hashrate distribution
\cite{milad2024statistical}. Bootstrap confidence intervals were computed for
every scenario, metric and monthly window (1{,}000 iterations), under two
resampling models as a cross-check: multinomial (fixed window total, the
exact block-level bootstrap for count-based shares) and Poisson (independent
per-actor counts). The two models' CI half-widths differ by 3.2\% on average,
confirming the intervals measure sampling variability rather than a resampling
artifact.
 
The intervals themselves tell a simple story (Figure~\ref{fig:ci}): they are
wide in the early era, when a month contained few blocks and coverage was poor,
and tight in the modern era, the baseline HHI carries a 95\% CI half-width
of roughly 0.004 from 2018 onwards, an order of magnitude smaller than even the
\emph{smallest} attribution band of \S\ref{sec:res-rq2}. Sampling and
attribution uncertainty do compound where coverage is low, the H2/H4
scenarios, especially under worst-case unknown assignment, run measurably
wider, but in the high-coverage era the effect is modest.
 
\begin{figure}[t]
  \centering
  \includegraphics[width=\linewidth]{bootstrap_ci_baseline.png}
  \caption{Baseline metrics with 95\% bootstrap sampling confidence intervals.
  Intervals are wide in the sparse early era and tight ($\mathrm{HHI}$
  half-width $\approx 0.004$) in the modern era.}
  \label{fig:ci}
\end{figure}
 
The same machinery closes RQ3, the pre-registered temporal-stability question.
Under the first criterion set, the weekly, monthly and quarterly granularities
pass while daily narrowly fails: window-to-window variation declines smoothly
with coarsening (HHI coefficient of variation 0.110 daily, just above the
strictest 0.10 ceiling, falling to 0.022 quarterly), rank correlations against
the monthly baseline are 0.94--0.998, and the 2014-11 and 2022-09 change-points are
detected identically at every granularity (figure in
Appendix~\ref{app:supp}). Under the proposal's bootstrap criterion,
adjacent-window 95\% CI overlap above 50\%, the monthly baseline passes for
every metric: Nakamoto 82.8\%, $\mathrm{CR}_3$ 61.2\%, $\mathrm{CR}_5$ 60.3\%,
HHI 56.9\%, Shannon 56.0\%. The single failure is Gini at 33.5\%, consistent
with its fragility throughout this chapter. Monthly aggregation is therefore
stable as pre-registered, and granularity choice is ruled out as a driver of
the trajectory.
 
%----------------------------------------------------------------------
\section{Robustness of structural breaks (RQ4)}
\label{sec:res-rq4}
 
The capstone question: are the turning points of \S\ref{sec:res-trajectory}
real features of Bitcoin's history, or artifacts of attribution? Because much
of the signal is smooth drift rather than sharp shift, demanding that detected
break \emph{dates} match exactly across scenarios would be an unfairly brittle
test; robustness is instead assessed by two complementary measures across all
fifteen attribution scenarios, on the high-coverage era that
\S\ref{sec:res-coverage} and RQ3 independently established as measurable.
 
\textbf{Trajectory agreement.} Rank correlation of each scenario's metric
series against the baseline is high where it should be: median Spearman $\rho =
0.912$ for HHI and $0.918$ for Shannon entropy across scenarios
(Figure~\ref{fig:changepoint}). The concentration ratios show low full-history
correlations ($\rho \approx 0.30$), but this is a saturation artifact rather
than disagreement: $\mathrm{CR}_3$ and $\mathrm{CR}_5$ pin near 1.0 in the
early era, destroying rank variation; restricting to 2014 onwards lifts them
into agreement with the rest. In every metric, the lowest-correlating scenarios
are consistently the low-coverage H2/H4 arms.
 
\textbf{Corroboration of the re-concentration break.} For the salient recent
break, scenario-by-scenario corroboration within a fixed tolerance window,
against a chance rate of roughly 13\%, reaches 14 of 15 scenarios for HHI
and 12 of 15 for $\mathrm{CR}_3$, far above chance. Shannon entropy and
$\mathrm{CR}_5$ corroborate in only 6 of 15, not because the shift is absent
but because it is smoother in those metrics, blurring the detected month.
 
\begin{figure}[t]
  \centering
  \includegraphics[width=\linewidth]{changepoint_robustness.png}
  \caption{Decentralisation trajectory across the fifteen attribution
  scenarios. Tag-based scenarios agree on the shape; divergence concentrates in
  the address-heavy, low-coverage scenarios.}
  \label{fig:changepoint}
\end{figure}
 
Two further observations locate the sensitivity precisely. Tag-only
attribution (H1) is the most stable heuristic, it never appears among the
bottom-four correlating scenarios for any metric, so the instability is
concentrated in the \emph{address signal}. And a deliberately unrestricted
all-history run confirms the scope restriction was necessary rather than
convenient: in the pre-2014 low-coverage era, detected break dates scatter by
years across scenarios, i.e.\ there is no robust break structure there to find.
 
The headline: \emph{the recent re-concentration is a real structural shift
whose trajectory is robust to attribution choice; the sensitivity that exists
is localised (the address signal, the low-coverage era); and the exact dating
of gradual transitions is attribution-dependent}. The strength of the evidence
varies by metric --- the break is corroborated in 14 of 15 scenarios for HHI
but, because it is smoother there, only 6 of 15 for Shannon entropy and
$\mathrm{CR}_5$ --- so the claim is properly read as trajectory-level agreement
that the shift is real and recent, not as month-exact agreement on its onset.
The analysis distinguishes which claims the data supports from which it cannot.
 
%----------------------------------------------------------------------
\section{Synthesis: an ordering of uncertainty}
\label{sec:res-synthesis}
 
The chapter's results assemble into a single comparison. Every uncertainty
source in the measurement --- the attribution heuristic, the unknown-block
treatment, the actor definition, and sampling noise --- has now been quantified
on the same data, the same metric, and the same period.
Table~\ref{tab:ordering} states the result for HHI over the high-coverage era.
 
\begin{table}[t]
  \centering
  \caption{The ordering of uncertainty sources (HHI, monthly windows, 2018
  onwards). Band widths are mean max--min across the relevant scenario axis;
  sampling is the mean 95\% bootstrap CI half-width.}
  \label{tab:ordering}
  \begin{tabular}{llr}
    \toprule
    Rank & Source of uncertainty & Magnitude (HHI) \\
    \midrule
    1 & Attribution heuristic (H1--H5)      & 0.272 \\
    2 & Unknown-block treatment (U1--U3)    & 0.037 \\
    3 & Actor definition (A1--A2 sweep)     & 0.015 \\
    4 & Sampling (bootstrap 95\% CI)        & $\approx$0.004 \\
    \bottomrule
  \end{tabular}
\end{table}
 
The ordering spans nearly two orders of magnitude:
\begin{equation}
  \underbrace{0.272}_{\text{heuristic}} \;\gg\;
  \underbrace{0.037}_{\text{unknown}} \;>\;
  \underbrace{0.015}_{\text{actor}} \;>\;
  \underbrace{0.004}_{\text{sampling}} .
  \label{eq:ordering}
\end{equation}
 
This inequality is the thesis in one line. \emph{How} blocks are attributed
matters one to two orders of magnitude more than how many blocks are sampled,
so a study that reports statistical confidence intervals while leaving its
attribution method undisclosed has quantified its smallest uncertainty and
hidden its largest. Within attribution, the method combination (which signals,
at what coverage) dominates unknown handling, which dominates actor definition.
The literature's sampling-uncertainty framework \cite{milad2024statistical} and
this thesis's attribution sweep are thus complementary halves of a complete
uncertainty account, and the ordering between them is not close.
Chapter~\ref{ch:discussion} translates this ordering into concrete reporting
guidance.


%======================================================================
\chapter{Discussion}
\label{ch:discussion}
% BUDGET ~5 pp.

The results of Chapter~\ref{ch:results} support one headline recommendation:
\emph{decentralisation studies should report sensitivity envelopes and
confidence intervals, not point estimates, and should disclose the attribution
heuristic used and the coverage it achieved.} A bare claim of the form
``Bitcoin's Nakamoto coefficient is 2'' is not a defensible measurement: as
this thesis shows, the same blockchain, on the same day, supports materially
different values of that number depending on choices the claim does not
mention. The recommendation is not a counsel of general caution, it follows
quantitatively from the uncertainty ordering of \S\ref{sec:res-synthesis}, and
the ordering also says \emph{which} disclosures matter most. Four specific
recommendations follow.

\textbf{Disclose the heuristic first.} Because the tag/address method
combination dominates every other uncertainty source, an HHI band of 0.272
against 0.037 for unknown handling, 0.015 for actor definition and
$\approx$0.004 for sampling (Table~\ref{tab:ordering}), the single most
consequential disclosure a study can make is how the two attribution signals
were consulted and combined. It is also the disclosure the literature most
conspicuously omits: studies routinely name their data source and their
metric, but almost never state whether attribution rested on coinbase tags,
payout addresses, or some resolution of the two. On the evidence of this
thesis, that omission hides the largest term in the error budget while
carefully reporting the smallest.

\textbf{Coverage must travel with every metric.} A share distribution computed
over an attributed subset is only informative about the whole population if
the subset is representative, and \S\ref{sec:res-rq2} shows it often is not:
the low-coverage heuristics (35.6\% and 29.2\% of blocks) do not merely add
noise around the high-coverage trajectory, they trace a systematically
different one. Coverage is thus the single best predictor of where in the
sensitivity envelope an estimate lies, and a reported metric without its
coverage figure cannot be interpreted, readers cannot tell whether they are
looking at a measurement of Bitcoin or a measurement of the attributable
third of it.

\textbf{Gini requires a specific warning.} The Gini coefficient is uniquely
exposed to the actor definition: a 37\% band under the A2 threshold sweep
against 4--9\% for the head-weighted metrics (\S\ref{sec:res-rq1}), and the
only metric to fail the pre-registered stability criterion
(\S\ref{sec:res-sampling}). The mechanism is structural, not incidental,
Gini is constructed from exactly the long-tail mass that actor and threshold
choices redistribute, so two researchers with identical data and different
long-tail handling will publish very different Gini values, both ``correct''.
Gini should not be reported for mining decentralisation without an explicit
statement of the actor definition and any inclusion threshold; where the
security-relevant head of the distribution is the object of interest, the
head-weighted metrics are simply better instruments.

\textbf{Report threshold metrics as ranges.} The Nakamoto coefficient's
integer, threshold-crossing construction makes it discontinuously sensitive to
choices that move mass near the 50\% boundary: its band under unknown-handling
alone reaches a full unit (\S\ref{sec:res-rq2}). This complements Milad et
al.'s statistical argument for reporting the Nakamoto coefficient as a
supported range rather than a point \cite{milad2024statistical}; the present
results extend that argument from sampling uncertainty to attribution
uncertainty, where the effect is larger.

There is also a positive reframing in these results. Sampling uncertainty,
the one uncertainty source the field has begun to formalise
\cite{milad2024statistical}, turns out to be small and well-behaved in the
modern era: tight intervals, consistent across resampling models, negligible
against attribution effects. The community's marginal effort is therefore
misallocated if it goes to larger samples or longer windows; the binding
constraint on measurement quality is attribution transparency, and that is
where effort should go. This is, in the end, an optimistic conclusion:
the dominant uncertainty is not an irreducible fact of nature but a reporting
practice, and reporting practices can change.

\section{Limitations}
\label{sec:disc-limitations}

\textbf{Operator clustering is confounded, a finding, not only a
limitation.} The most important limitation of this work is that it cannot
deliver a clean operator-level (A3) measurement, and the reason is itself one
of the thesis's results. Linking pool labels that share payout addresses,
the natural operationalisation of the SoK's open operator-identity question
\cite{ovezik2025sok}, systematically over-merges independent major pools:
connected components fused 64\% of all blocks into a single implausible
``operator'' spanning AntPool, F2Pool, ViaBTC, Binance and others, and Louvain
community detection \cite{blondel2008louvain} reduced but did not eliminate the
fusion (47\% of blocks). Nor can the problem be tuned away: requiring at least
2, 3 or 5 shared addresses per link shrank the largest cluster from 28 labels
to 19, 17 and 12, but never broke the core, because the edges joining the
majors carry weights of 1{,}145--1{,}556 shared addresses. These are not weak
bridges to be thresholded away; major pools genuinely share
hundreds-to-thousands of payout addresses on-chain, through shared custodial
and mining-as-a-service infrastructure and address reuse across rebrands. The
signal is confounded at the top of the distribution, precisely where the
decentralisation question lives, rather than noisy, so no parameter choice
fixes it. This is consistent with independent findings on address clustering
more broadly: Harrigan and Fretter~\cite{harrigan2016unreasonable} document how
address reuse and ``avoidable merging'' produce high-centrality super-clusters,
and observe that a merge of two already-large clusters is itself a signature of
a spurious rather than genuine link --- exactly the failure mode seen here when
independent major pools are fused. What the method \emph{does} recover, consistently across both
clustering algorithms, are small and historically plausible merges such as
\{Eligius, KnCMiner\} and \{NiceHash, Solo~CK\}, which is exactly the
behaviour of a signal that is informative in the tail and confounded at the
head. The mitigation is already built into the thesis: A3 results are reported
strictly as a sensitivity upper bound, never as a measurement, and the clean
answer to RQ1 rests on A1 and A2 (\S\ref{sec:res-rq1}). Framed positively,
this is an empirical answer to a question the SoK leaves open: the natural
method was tested at full-history scale and shown, with evidence, where and
why it fails, credible operator attribution needs evidence beyond shared
addresses alone, and Chapter~\ref{ch:conclusions} specifies what that evidence
should be.

\textbf{Reliance on public reference lists with shared lineage.} All
attribution in this thesis descends from two public reference lists, and those
lists share community lineage, the mempool.space list is a descendant of
the same open-source effort as the btc.com list. This is a genuine dependence:
the lists are two snapshots of one curation ecosystem, not two independent
instruments, and the near-perfect agreement found in \S\ref{sec:res-interlist}
is in part explained by exactly this shared ancestry. The thesis handles the
dependence by naming it rather than hiding it: the inter-list result is
framed as measuring the convergence of the public attribution ecosystem
(against Romiti et al.'s 684 conflicts across the older, heterogeneous source
landscape \cite{romiti2019deep}), not the concordance of independent
observers. The practical conclusion survives the caveat, whatever
uncertainty remains in attribution, it demonstrably does not live in the
choice between today's major public lists.

\textbf{The early era is unreliable, and the conclusions say so.} Before
roughly 2014, coverage is structurally poor, pools were few and
loosely-labelled, and \S\ref{sec:res-rq4} shows that detected break dates
there scatter by years across attribution scenarios, there is no robust
structure to recover. Rather than truncating the data silently, the full
history is processed, the coverage series is reported, and every strong claim
in this thesis is explicitly scoped to the high-coverage era. The limitation
is thus converted into a delimitation: the thesis states where measurement is
supportable and declines to claim beyond it.

\section{Threats to validity}
\label{sec:disc-validity}

\textbf{Reference-list completeness.} Thirty percent of all blocks in history
match no tag and no known address, and the measured picture is conditional on
what the lists can see. If unattributed blocks were produced
disproportionately by unlisted large actors, the true concentration could
differ from the baseline estimate. This threat cannot be eliminated, it is
the fundamental epistemic condition of the problem, but it is bounded
rather than ignored: the U1--U3 sweep propagates the unknown mass through to
the metrics under assumptions spanning anonymous-entity, missing-at-random and
adversarial worst-case readings, so the reported envelope already contains
the range of pictures consistent with the unattributed mass. In the modern,
high-coverage era the envelope is narrow; the threat is material exactly
where the thesis already declines to make strong claims.

\textbf{Smooth signals resist exact dating.} The change-point analysis can
robustly establish \emph{that} the recent re-concentration is a real
structural shift, trajectory agreement and above-chance corroboration
across attribution scenarios (\S\ref{sec:res-rq4}), but it cannot pin the
shift to an exact month, because a gradual transition has no crisp date to
find: on smooth drift, the detected boundary moves with minute perturbations
of the series. This is an honest epistemic boundary rather than a
methodological defect, and the method was designed around it
(\S\ref{sec:changepoint}): RQ4 deliberately evaluates trajectory agreement
and tolerance-window corroboration instead of exact date matching, so the
robustness claim made is one the data can actually support. The conclusion is
stated at the precision the signal permits, the re-concentration is real;
its exact onset month is attribution-sensitive, and no more.

\textbf{Scope of the conclusions.} The strong conclusions of this thesis,
the U-shaped trajectory, the reality of the re-concentration, the uncertainty
ordering, hold for pool-level attribution, in the high-coverage era, at
monthly granularity. They are not claims about the pre-2014 era (where the
thesis shows measurement is unsupportable), about operator-level concentration
(where the thesis shows the natural method is confounded and reports only
bounds), or about the intra-pool distribution of hashers (invisible on-chain,
\S\ref{sec:data}). Each scope restriction is not an assumption but a result:
in every case the boundary is drawn where the thesis's own analysis
demonstrated measurement breaks down. This is the discipline the thesis
recommends to the field, applied to itself, claims accompanied by the
evidence of where they stop.


%======================================================================
\chapter{Conclusions}
\label{ch:conclusions}
% Drop-in replacement for the Conclusions chapter in skeleton.tex.
% BUDGET ~3 pp. Summary of contributions -> future work (HR4/HR5 framed as
% principled deferrals building on stronger operator evidence).

\section{Summary of contributions}
\label{sec:conc-summary}

This dissertation set out from a simple observation: every claim about the
decentralisation of Bitcoin mining rests on attributing each block to a pool,
and the literature treats that step as settled when it is not. By making
attribution the object of study, varying it systematically along three
dimensions over the full block history and propagating the variation through
to the final metrics, the thesis delivers five contributions.

\textbf{An empirical picture with its trust boundary attached.} Bitcoin mining
follows a U-shaped trajectory: decentralising from its early years to its most
decentralised state around 2017, and re-concentrating since, to a present state
near a Nakamoto coefficient of 2 with the top three pools producing roughly 70\% of blocks.
Crucially, the analysis states where this picture is reliable and where it is
not: the trajectory and its turning points are robust to attribution choice in
the high-coverage era, the recent re-concentration is corroborated far above
chance across attribution scenarios, and the pre-2014 era admits no stable
conclusions at all.

\textbf{A quantitative ordering of uncertainty.} The central result is that
every source of uncertainty in the measurement was quantified on the same
data and can therefore be ranked. For HHI over the high-coverage era: the
attribution heuristic contributes a band of 0.272, unknown-block handling
0.037, actor definition 0.015, and sampling noise approximately 0.004
(Equation~\ref{eq:ordering}). How blocks are attributed matters one to two
orders of magnitude more than how many blocks are sampled, and within
attribution the tag/address method combination dominates everything else.
This complements the sampling-uncertainty framework of Milad et
al.~\cite{milad2024statistical}: the two analyses quantify independent
uncertainty sources, and this thesis establishes the ordering between them.

\textbf{The convergence of public attribution sources.} Against Romiti et
al.'s historical finding of 684 inter-source conflicts
\cite{romiti2019deep}, the two major public reference lists of today agree on
99.9997\% of jointly covered blocks, two conflicts in all of history after
resolving one rebrand alias. Attribution uncertainty is real, but it no
longer lives in list choice; it lives in the tag-versus-address split,
coverage, and unknown handling. The shared lineage of the lists is named as
both the caveat and the explanation.

\textbf{An empirical answer to the operator-identity question.} The SoK
framework names entity clustering as the open challenge of actor
identification \cite{ovezik2025sok}. This thesis tested the natural method,
clustering pool labels by shared payout addresses, at full-history scale
and showed, with evidence, where and why it fails: independent major pools
genuinely share hundreds to thousands of payout addresses through custodial
and mining-as-a-service infrastructure, so the method over-merges precisely at
the top of the distribution and no threshold repairs it. Operator-level
results are therefore reported as bounds, and the finding stands as a negative
result the field can build on: credible operator attribution requires
evidence beyond shared addresses alone.

\textbf{Reporting guidance grounded in measurement.} The practical output is
a set of recommendations that follow from the ordering rather than from
general caution: report sensitivity envelopes and confidence intervals, never
bare point estimates; disclose the attribution heuristic first, because it is
the largest term in the error budget; attach coverage to every metric; treat
Gini as unsafe without an explicit actor definition; and report threshold
metrics like the Nakamoto coefficient as supported ranges. All code, derived
data, and pre-registered thresholds are publicly archived, so the thesis
applies its own standard: every number it reports can be traced to a
disclosed configuration and reproduced.

\section{Future work}
\label{sec:conc-future}

Two extensions were identified at proposal stage and deliberately deferred;
the results of this thesis now say precisely what they require, which is what
makes them principled next steps rather than omissions.

\textbf{A temporal operator graph.} The operator-clustering analysis treated
address sharing as a static, all-history signal, and the confound it uncovered
is partly a temporal artifact: addresses inherited across rebrands and
custodial infrastructure adopted at different times are flattened into a
single graph. A time-sliced address-sharing graph, tracking when sharing
relationships form and dissolve, would separate genuine contemporaneous
operator overlap from historical inheritance, and would turn the
shared-infrastructure phenomenon itself, documented here as a confound, into
an object of study. The small, historically plausible merges that both
clustering methods recovered provide a natural validation set.

\textbf{Probabilistic attribution.} The deterministic heuristic sweep was the
right instrument for this thesis's question, and the reason for deferring its
probabilistic generalisation, that an inferential model layers new,
hard-to-validate modelling uncertainty on top of the uncertainty under study,
is now backed by evidence: the dominant uncertainty is the method
combination itself, which a probabilistic model must weight somehow. A
credible probabilistic attribution therefore needs what this thesis showed the
on-chain signals alone cannot supply: operator-level evidence beyond shared
addresses, such as curated off-chain records of rebrands and corporate
relationships, combined with the temporal graph above. Given such evidence,
each block could carry a distribution over producers, and the sensitivity
envelopes of this thesis would become posterior credible intervals, with
the confidence-interval machinery of Chapter~\ref{ch:results} extending
naturally to that setting.

Beyond these two, the framework itself travels. The three-dimensional design, 
who counts as an actor, which signals attribute, what happens to the
unattributed, is not specific to Bitcoin, and applying it to other
proof-of-work chains, or to other resources in the stratified sense of
Karakostas et al.~\cite{ovezik2022sok}, would test whether the uncertainty
ordering found here is a fact about Bitcoin or a fact about blockchain
measurement in general. A further watch item is structural: emerging
template-negotiation protocols such as Stratum~v2 may decouple block
construction from pool identity, degrading tag-based attribution, currently
the most stable signal, and shifting the balance of evidence toward exactly
the operator-level methods this section calls for.

The closing point is the thesis in one sentence: decentralisation is
measurable, but only honestly, with the attribution method disclosed, the
coverage attached, and the uncertainty reported, and the largest share of
that uncertainty is, encouragingly, within the community's power to remove.


%======================================================================
\bibliographystyle{plain}
\bibliography{mybibfile}


% You may delete everything from \appendix up to \end{document} if you don't need it.
\appendix

\chapter{Supplementary material}
\label{app:supp}
% Optional. Markers are not obliged to read appendices, so keep contributions in the main
% body. Candidates for here: the BigQuery partition-pruning detail, the full 15-scenario
% envelope figure (sensitivity_full_envelope.png), the granularity-stability figure
% (stability_granularity.png), the unknown-only sensitivity figure (sensitivity_unknown.png),
% and the A1/A2/A3 full actor band (actor_definition_full_band.png, labelled as an upper bound).

% TODO(you): Appendix A is currently empty but is cited from the main body
% (\S\ref{sec:res-sampling} refers to "figure in Appendix~\ref{app:supp}", and
% \S\ref{sec:data} to "details in Appendix~\ref{app:supp}"). Either populate it
% with the figures/detail below or remove those forward-references. A stub
% figure environment is provided so the cross-reference resolves; replace the
% filename with your real granularity-stability figure.

\section{Temporal-stability across granularities}
\label{app:stability}
% TODO(you): drop in stability_granularity.png (daily/weekly/monthly/quarterly
% metric series with the 2014-11 and 2022-09 change-points marked).
\begin{figure}[h]
  \centering
  % \includegraphics[width=\linewidth]{stability_granularity.png}
  \caption{Metric trajectories at daily, weekly, monthly and quarterly
  granularity, with detected change-points marked. Referenced from
  \S\ref{sec:res-sampling}. \textbf{TODO: insert figure file.}}
  \label{fig:stability-granularity}
\end{figure}

\section{Data and reproducibility}
\label{app:repro}
% TODO(you): add the BigQuery partition-pruning / extraction-cost detail cited
% from \S\ref{sec:data}, and the archive location. The main body repeatedly
% states that "code, derived data, and pre-registered thresholds are publicly
% archived" (\S\ref{sec:pipeline}, \S\ref{sec:conc-summary}); that claim needs a
% resolvable link/DOI to be assessable.
%
% All code, derived datasets, and the timestamped pre-registered thresholds are
% available at: \url{TODO-INSERT-REPOSITORY-OR-DATASHARE-DOI}

\end{document}
